%No specific subdivision. 
%Key summary on introduction (problem statement), research objectives, methodology, results, and conclusions.

%\DD{}{Maybe add a short sentence introducing churn and the importance of predicting it? Check the limit in terms of words} -Closed
\noindent To retain customers effectively, retailers need accurate churn prediction to drive personalized engagement strategies. This study presents a systematic literature review of artificial intelligence applications for customer churn prediction in the retail sector, with a focus on non-contractual settings where churn must be detected behaviorally rather than explicitly recorded. From an initial pool of 1,370 articles identified across four major scientific databases, following a rigorous selection process, 41 relevant studies were selected and analyzed to answer \textcolor{blue}{8} research questions drawn from our research goal. The review synthesizes findings on churn definitions, datasets, feature engineering practices, AI techniques, and evaluation metrics used for churn prediction. A notable trend is the reliance on static churn definitions based on fixed inactivity time windows, despite the large use of dynamic, personalized approaches that more accurately capture customer behavior. The widespread use of feature engineering methods, such as feature transformation, demonstrates the potential of this process to enhance model performance.
Tree-based models and neural network architectures are identified as the most effective predictive techniques. However, heterogeneity in datasets, churn definitions, and modeling approaches poses challenges for performance comparisons. Additional issues include limited access to real-world datasets, class imbalance, and a lack of explainable AI tools. The research in this field is evolving to address these limitations with more sophisticated approaches. Future research directions proposed by authors include expanding model diversity (e.g., employing transformers, RNNs, and ensembles), improving explainability, and developing automated frameworks for feature selection and hyperparameter tuning.







%This review offers a comprehensive overview of the current landscape and proposes future research directions, including the standardization of churn labeling, integration of contextual data, and greater emphasis on model interpretability. It aims to support both academic inquiry and practical implementation efforts aimed at improving customer retention in increasingly competitive and data-driven retail environments.



%no more than 250 words.  It should briefly state the purpose of the research, principal results, and major conclusions.

%Keywords: Include 1 to 7 keywords for indexing purposes. Avoid multi-word phrases and non-standard abbreviations.


